% GENERATED FILE. Edit canonical lessons, then run npm run generate:books.
% canonical-source-hash: fnv1a64:cb77fb5b9d122f6a
% canonical-lessons: RU-C399-set, RU-C399-sayt, RU-C399-ssylka, RU-C399-brauzer, RU-C399-poiskovik

\chapter{Online}
\label{ch:ru-online}
\hlchaptermodality{399}

\begin{chapteropening}
\textbf{By the end of this chapter:} \emph{I can talk about the internet, a website and a link, a browser and a search engine.}

Complete the last lesson of chapter 399: I can talk about the internet, a website and a link, a browser and a search engine.
\end{chapteropening}

\section[\texorpdfstring{\ru{сеть} (set)}{сеть (set)}]{\ru{сеть} (set) — a network; the internet; a net}

\label{lesson:RU-C399-set}



\begin{quote}
Before the new one: say the Russian for to insure, then the Russian for insurance.
\end{quote}

\subsection*{You'll want to know: \ru{сеть}}
\textbf{\ru{сеть}} (\emph{set}) — \textquotedblleft{}a network; the internet; a net\textquotedblright{}.

\begin{quote}
\ru{Я} \ru{нашёл} \ru{эту} \ru{информацию} \ru{в} \ru{сети}.
\end{quote}

\begin{quote}
\emph{Ya nashyól étu informátsiyu v setí.}
\end{quote}

\begin{quote}
I found this information online.
\end{quote}

Feminine. \emph{\ru{В} \ru{сети}} — online. \emph{\ru{Социальная} \ru{сеть}} — a social network. First of all a fishing net.

\begin{cousinweb}
\textbf{\ru{сеть}} is inherited from Proto-Slavic \emph{*sětĭ} \textquotedblleft{}a net, a snare\textquotedblright{} (EtymDB). The internet is a net in the old fisherman’s sense.
\end{cousinweb}

\begin{grammarlens}[title={What you've built}]
One more word for telling, explaining and describing at B1.
\end{grammarlens}

\subsection*{Your turn}
\begin{itemize}
\raggedright
  \item \emph{Say it:} \emph{set}
  \item \emph{Say it:} \emph{set}, once more
  \item \emph{Say it:} a sentence of your own with \emph{set}
  \item \emph{Recall:} say the Russian for to insure, then the Russian for insurance, then say \emph{set} again
\end{itemize}

\begin{tcolorbox}[breakable,colback=teal!4,colframe=teal!35!black,title={Before you move on}]
What does \ru{сеть} mean? (\textquotedblleft{}A network; the internet; a net\textquotedblright{}.) Say it once more.
\end{tcolorbox}
\section[\texorpdfstring{\ru{сайт} (sayt)}{сайт (sayt)}]{\ru{сайт} (sayt) — a website}

\label{lesson:RU-C399-sayt}



\begin{quote}
Before the new one: say the Russian for insurance, then the Russian for a network.
\end{quote}

\subsection*{You'll want to know: \ru{сайт}}
\textbf{\ru{сайт}} (\emph{sayt}) — \textquotedblleft{}a website\textquotedblright{}.

\begin{quote}
\ru{У} \ru{нашей} \ru{школы} \ru{есть} \ru{свой} \ru{сайт}.
\end{quote}

\begin{quote}
\emph{U náshey shkóly yest svoy sayt.}
\end{quote}

\begin{quote}
Our school has its own website.
\end{quote}

Masculine. \emph{\ru{На} \ru{сайте}} — on the website.

\begin{cousinweb}
\textbf{\ru{сайт}} was borrowed from English \emph{site} (EtymDB), which is usually traced, through French, to Latin \emph{situs} \textquotedblleft{}a place, a position\textquotedblright{} (Lewis \& Short).
\end{cousinweb}

\begin{grammarlens}[title={What you've built}]
One more word for telling, explaining and describing at B1.
\end{grammarlens}

\subsection*{Your turn}
\begin{itemize}
\raggedright
  \item \emph{Say it:} \emph{sayt}
  \item \emph{Say it:} \emph{sayt}, once more
  \item \emph{Say it:} a sentence of your own with \emph{sayt}
  \item \emph{Recall:} say the Russian for insurance, then the Russian for a network, then say \emph{sayt} again
\end{itemize}

\begin{tcolorbox}[breakable,colback=teal!4,colframe=teal!35!black,title={Before you move on}]
What does \ru{сайт} mean? (\textquotedblleft{}A website\textquotedblright{}.) Say it once more.
\end{tcolorbox}
\section[\texorpdfstring{\ru{ссылка} (ssýlka)}{ссылка (ssýlka)}]{\ru{ссылка} (ssýlka) — a link; a reference; exile}

\label{lesson:RU-C399-ssylka}



\begin{quote}
Before the new one: say the Russian for a network, then the Russian for a website.
\end{quote}

\subsection*{You'll want to know: \ru{ссылка}}
\textbf{\ru{ссылка}} (\emph{ssýlka}) — \textquotedblleft{}a link; a reference; exile\textquotedblright{}.

\begin{quote}
\ru{Я} \ru{отправил} \ru{тебе} \ru{ссылку}.
\end{quote}

\begin{quote}
\emph{Ya otprávil tebé ssýlku.}
\end{quote}

\begin{quote}
I have sent you the link.
\end{quote}

Feminine. \emph{\ru{Ссылка} \ru{на}} + accusative — a link to, a reference to. Historically, exile: \emph{\ru{отправить} \ru{в} \ru{ссылку}}.

\begin{cousinweb}
\textbf{\ru{ссылка}} is built on \emph{\ru{ссылать}} \textquotedblleft{}to send away\textquotedblright{} and \emph{\ru{ссылаться}} \textquotedblleft{}to refer\textquotedblright{}, from \emph{\ru{слать}} \textquotedblleft{}to send\textquotedblright{}, the root of \emph{\ru{посылка}}. A link sends you on to another page.
\end{cousinweb}

\begin{grammarlens}[title={What you've built}]
One more word for telling, explaining and describing at B1.
\end{grammarlens}

\subsection*{Your turn}
\begin{itemize}
\raggedright
  \item \emph{Say it:} \emph{ssýlka}
  \item \emph{Say it:} \emph{ssýlka}, once more
  \item \emph{Say it:} a sentence of your own with \emph{ssýlka}
  \item \emph{Recall:} say the Russian for a network, then the Russian for a website, then say \emph{ssýlka} again
\end{itemize}

\begin{tcolorbox}[breakable,colback=teal!4,colframe=teal!35!black,title={Before you move on}]
What does \ru{ссылка} mean? (\textquotedblleft{}A link; a reference; exile\textquotedblright{}.) Say it once more.
\end{tcolorbox}
\section[\texorpdfstring{\ru{браузер} (bráuzer)}{браузер (bráuzer)}]{\ru{браузер} (bráuzer) — a web browser}

\label{lesson:RU-C399-brauzer}



\begin{quote}
Before the new one: say the Russian for a website, then the Russian for a link.
\end{quote}

\subsection*{You'll want to know: \ru{браузер}}
\textbf{\ru{браузер}} (\emph{bráuzer}) — \textquotedblleft{}a web browser\textquotedblright{}.

\begin{quote}
\ru{Этот} \ru{сайт} \ru{не} \ru{открывается} \ru{в} \ru{моём} \ru{браузере}.
\end{quote}

\begin{quote}
\emph{Étot sayt ne otkryváyetsya v moyóm bráuzere.}
\end{quote}

\begin{quote}
This website does not open in my browser.
\end{quote}

Masculine. From English \emph{browser}.

\begin{cousinweb}
\textbf{\ru{браузер}} was borrowed from English \emph{browser} (EtymDB) in the 1990s, keeping the English sound.
\end{cousinweb}

\begin{grammarlens}[title={What you've built}]
One more word for telling, explaining and describing at B1.
\end{grammarlens}

\subsection*{Your turn}
\begin{itemize}
\raggedright
  \item \emph{Say it:} \emph{bráuzer}
  \item \emph{Say it:} \emph{bráuzer}, once more
  \item \emph{Say it:} a sentence of your own with \emph{bráuzer}
  \item \emph{Recall:} say the Russian for a website, then the Russian for a link, then say \emph{bráuzer} again
\end{itemize}

\begin{tcolorbox}[breakable,colback=teal!4,colframe=teal!35!black,title={Before you move on}]
What does \ru{браузер} mean? (\textquotedblleft{}A web browser\textquotedblright{}.) Say it once more.
\end{tcolorbox}
\section[\texorpdfstring{\ru{поисковик} (poiskovík)}{поисковик (poiskovík)}]{\ru{поисковик} (poiskovík) — a search engine}

\label{lesson:RU-C399-poiskovik}



\begin{quote}
Before the new one: say the Russian for a link, then the Russian for a web browser.
\end{quote}

\subsection*{You'll want to know: \ru{поисковик}}
\textbf{\ru{поисковик}} (\emph{poiskovík}) — \textquotedblleft{}a search engine\textquotedblright{}.

\begin{quote}
\ru{Я} \ru{ищу} \ru{ответ} \ru{в} \ru{поисковике}.
\end{quote}

\begin{quote}
\emph{Ya ishchú otvét v poiskovíke.}
\end{quote}

\begin{quote}
I am looking for the answer in a search engine.
\end{quote}

Masculine, colloquial. Built on \emph{\ru{поиск}}, a search.

\begin{cousinweb}
\textbf{\ru{поисковик}} is built on \emph{\ru{поисковый}} \textquotedblleft{}search\textquotedblright{}, from \emph{\ru{поиск}} \textquotedblleft{}a search\textquotedblright{}, from \emph{\ru{искать}} \textquotedblleft{}to look for\textquotedblright{}, inherited from Proto-Slavic \emph{*jĭskati} \textquotedblleft{}to look for\textquotedblright{} (EtymDB).
\end{cousinweb}

\begin{grammarlens}[title={What you've built}]
One more word for telling, explaining and describing at B1.
\end{grammarlens}

\subsection*{Your turn}
\begin{itemize}
\raggedright
  \item \emph{Say it:} \emph{poiskovík}
  \item \emph{Say it:} \emph{poiskovík}, once more
  \item \emph{Say it:} a sentence of your own with \emph{poiskovík}
  \item \emph{Recall:} say the Russian for a link, then the Russian for a web browser, then say \emph{poiskovík} again
\end{itemize}

\begin{tcolorbox}[breakable,colback=teal!4,colframe=teal!35!black,title={Before you move on}]
What does \ru{поисковик} mean? (\textquotedblleft{}A search engine\textquotedblright{}.) Say it once more.
\end{tcolorbox}
